\documentclass[12pt,letterpaper]{book}

\usepackage[utf8]{inputenc}
\usepackage[T1]{fontenc}
\usepackage{lmodern}
\usepackage[margin=1.25in,headheight=15pt]{geometry}
\usepackage{amsmath,amsthm,amssymb}
\usepackage{enumitem}
\usepackage{hyperref}
\usepackage{fancyhdr}
\usepackage{datetime}
\usepackage{csquotes}
\usepackage{titlesec}
\usepackage{setspace}

% Spacing
\onehalfspacing

% Chapter and section formatting
\titleformat{\chapter}[display]
  {\normalfont\huge\bfseries\centering}
  {\chaptertitlename\ \thechapter}{20pt}{\Huge}
\titleformat{\section}
  {\normalfont\Large\bfseries\centering}
  {\thesection}{1em}{}
\titleformat{\subsection}
  {\normalfont\large\bfseries}
  {\thesubsection}{1em}{}

% Spacing around titles
\titlespacing*{\chapter}{0pt}{50pt}{40pt}
\titlespacing*{\section}{0pt}{30pt}{20pt}
\titlespacing*{\subsection}{0pt}{20pt}{10pt}

% Header and footer setup
\pagestyle{fancy}
\fancyhf{}
\fancyhead[LE,RO]{\thepage}
\fancyhead[RE]{\textit{The Myth of Sisyphus}}
\fancyhead[LO]{\textit{Andre Ruiz Loera}}
\renewcommand{\headrulewidth}{0.4pt}

% Theorem environments
\newtheorem{theorem}{Theorem}
\newtheorem{proposition}{Proposition}
\theoremstyle{definition}
\newtheorem{definition}{Definition}
\theoremstyle{remark}
\newtheorem{remark}{Remark}
\newtheorem{note}{Note}

\hypersetup{
    colorlinks=true,
    linkcolor=blue,
    filecolor=magenta,
    urlcolor=cyan,
    pdftitle={The Myth of Sisyphus - Commentary},
    pdfauthor={Andre Ruiz Loera},
}

\begin{document}

% Title page
\begin{titlepage}
\begin{center}

\vspace*{4cm}

{\Huge\bfseries Commentary on}\\[1cm]

{\Huge\bfseries\textit{The Myth of Sisyphus}}\\[0.5cm]

{\large\textit{by Albert Camus}}

\vfill

{\LARGE Andre Ruiz Loera}

\vspace{3cm}

\end{center}
\end{titlepage}

% Foreword (before TOC)
\chapter*{Foreword}
\addcontentsline{toc}{chapter}{Foreword}

Write about how Camus frames the essay as an investigation into the \enquote{absurd,} and how he rejects both religion and traditional philosophy as sufficient answers. Mention that he sets out to examine how a person should live once meaning is gone.

\clearpage

% Table of Contents
\tableofcontents
\clearpage

\chapter{An Absurd Reasoning}

\section{Absurdity and Suicide}

Explain that Camus begins with the central question: if life has no meaning, is suicide the logical response? Show how he argues that the absence of meaning is not a reason to die but a starting point for thinking clearly.

\section{The Absurd Wall}

Discuss how human beings desire clarity, order, and explanations, yet the world remains silent and indifferent. Camus calls this clash \enquote{the absurd,} and this section shows where that feeling comes from.

\section{Philosophical Suicide}

Write about Camus' critique of philosophers who escape the absurd with leaps of faith (like Kierkegaard or Jaspers). He thinks turning to God, transcendence, or hope is a way of avoiding reality — a form of \enquote{suicide} of the mind.

\section{The Absurd Freedom}

Explain Camus' conclusion that once we accept life has no ultimate meaning, we become truly free. Without hope, we gain a new kind of freedom, responsibility, and intensity in living.

\chapter{The Absurd Man}

\section{Don Juanism}

Write about Don Juan as an example of someone who lives fully without seeking lasting meaning — he seeks \enquote{quantity} of experiences, not \enquote{quality} or eternal commitments.

\section{Drama}

The actor becomes a symbol of the person who lives many lives but owns none permanently. This section shows how embracing impermanence aligns with absurd freedom.

\section{Conquest}

The conqueror (or rebel) represents action, risk, and engagement with the world despite its meaninglessness. Write how this figure embodies living intensely without illusions.

\chapter{Absurd Creation}

\section{Philosophy and Fiction}

Explain how art can express the absurd but cannot resolve it. Camus argues that creative work becomes meaningful precisely because it has no eternal purpose.

\section{Kirilov}

Write about Kirilov (from Dostoevsky's \textit{Demons}) as someone who tries to \enquote{become God} by committing suicide to prove absolute freedom — an extreme, tragic attempt to escape the absurd.

\section{Dostoevsky}

Discuss how Dostoevsky deeply understands the absurd but ultimately escapes into Christian faith, which Camus sees as another form of philosophical suicide.

\chapter{The Myth of Sisyphus}

Explain that Sisyphus, condemned to repeat a meaningless task forever, symbolizes the human condition. Camus concludes that the right response to absurdity is revolt, consciousness, and inner freedom — \enquote{One must imagine Sisyphus happy.}

\chapter{Personal Reflections}

% Your own thoughts and interpretations

\end{document}
